\begin{abstract}
In the context of the widespread application of deep learning in streamflow prediction, whether these models truly learn hydrologically meaningful information remains unclear. Existing studies largely rely on post hoc techniques such as SHAP and LIME, lacking systematic and verifiable interpretability analysis, and failing to assess whether model responses are physically consistent or whether the extracted knowledge is transferable across models.To address this gap, this study treats interpretability as a primary objective and proposes a Sparse Transformer-based Explainable Hydrology framework (STEH). The framework links inputs, internal circuit-level structures, and outputs through structural sparsification and pathway-level analysis, enabling interpretable characterization. It further incorporates circuit-level ablation, sensitivity analysis, and cross-model validation to evaluate the verifiability, behavioral consistency, and transferability of learned knowledge.Experiments across large-scale basins show that structurally simpler models exhibit clearer internal pathways and stronger physical interpretability, while more complex models may introduce redundancy and even degrade performance. Moreover, physically meaningful information is consistently identified across models with varying performance levels. The extracted structural knowledge can be transferred to improve traditional machine learning models.These findings suggest that deep learning can serve not only as a predictive tool but also as apotential source of hydrological insight, highlighting the need to move from descriptive explanation toward systematic verification in interpretability research.
\end{abstract}

\noindent\textbf{Keywords:} Model interpretability, Explainable artificial intelligence, Structural sparsity, Hydrological modeling, Streamflow prediction, Transformer